\section{Auxiliary probability estimates}\label{app:probability}

\begin{lemma}[Finite Gaussian association]\label{pr:association}
Let $Z$ be a centered nondegenerate Gaussian vector whose covariance
matrix $\Sigma$ has nonnegative entries. For bounded coordinatewise
nondecreasing Borel functions $f,g$,
\[
 \mathbb E[f(Z)g(Z)]\ge\mathbb Ef(Z)\,\mathbb Eg(Z).
\]
The assertion also holds whenever truncation makes both sides converge.
\end{lemma}
\begin{proof}
For smooth bounded functions with bounded derivatives, let $(X_t,Y_t)$
have diagonal covariance blocks $\Sigma$ and cross block $t\Sigma$,
$0\le t\le1$. This is a positive semidefinite covariance: one may take
$X_t=Z$ and $Y_t=tZ+\sqrt{1-t^2}Z'$ with $Z'$ an independent copy.
Gaussian integration by parts gives
\[
 \frac d{dt}\mathbb E[f(X_t)g(Y_t)]
 =\sum_{i,j}\Sigma_{ij}
       \mathbb E[\partial_i f(X_t)\partial_jg(Y_t)]\ge0.
\]
The two off-diagonal covariance blocks account for the absence of a
factor $1/2$. The covariance derivative identity follows by
differentiating $\exp(-\langle\xi,C(t)\xi\rangle/2)$ and integrating
twice by parts, first for smooth compactly supported test functions;
cutoff approximation and covariance regularization give the displayed
form. At $t=0$ the two vectors are independent; at $t=1$ they coincide.
Integration proves the claim for smooth functions.

Convolution with a positive smooth product kernel preserves
monotonicity. Bounded coordinatewise monotone functions have a
Lebesgue-null discontinuity set, so the convolutions converge almost
everywhere for the nonsingular Gaussian law. Dominated convergence
proves the Borel assertion. At distinct finite Ornstein--Uhlenbeck
times the covariance is positive definite; coincident times can simply
be deleted. Applying this finite result on nested dense time grids and
using path continuity proves its survival applications. This argument
does not assert association after arbitrary conditioning.
\end{proof}

\subsection{Joint exact-total comparisons and small tubes}

Let $f_{r,\tau}$ denote the density of a sum of $r$ independent
exponentials of mean $\tau$. If $n$ marks have exact low count $k$,
$\pi=k/n$, and their $n+1$ gaps have total $P=(n+1)\tau$, exposing
$m$ specified marks and $m$ distinct specified gaps has joint density
\begin{equation}
 \frac{f_{n+1-m,\tau}(P-T)}{f_{n+1,\tau}(P)}
 \frac{\mathbb P\{\operatorname{Bin}(n-m,\pi)=k-K\}}
      {\mathbb P\{\operatorname{Bin}(n,\pi)=k\}}
 \label{bridge:prefix-density}
\end{equation}
relative to independent reference gaps and marks. Here $T,K$ are the
exposed totals, and the ratio is zero outside its support. Independence
of the two unobserved sums proves the identity. The same formula
with $m$ replaced by $2m$ applies jointly to the first $m$ and last $m$
marks and to the distinct gap sets
\[
 E_1,\ldots,E_m;\qquad E_{n-m+2},\ldots,E_{n+1}.
\]
There are $n-2m$ residual marks and $n+1-2m$ residual gaps.

For $m\le n/4$, or for $m=\lfloor n/8\rfloor$ at each end, this
ratio is bounded above on its entire support, uniformly on compact
positive $\tau$ and interior $\pi$. On each fixed central box
\[
 |T-m\tau|+|K-m\pi|\le C\sqrt n
\]
it is bounded below positively for sufficiently large $n$; in the
two-end formula replace $m$ by $2m$. Indeed the Gamma maximum density
is $O((1+r)^{-1/2})$, its density at the reference mean is
$\asymp(1+r)^{-1/2}$, and the corresponding binomial bounds have
the same order. Central Stirling bounds give the lowers after the
remaining totals are inside their support. These denominators are
at their exact reference means. Near-central variants require the
large-total hypotheses of Lemma~\ref{pr:gamma}.

\begin{lemma}[Exact multicolor small tube]\label{bridge:small-tube}
Fix $r$, $\epsilon>0$, $0<\tau_-\le\tau_+<\infty$, and $a>0$.
Let a word of length $n$ be uniform of exact type
$v=(v_1,\ldots,v_r)$, $\pi_c=v_c/n\ge\epsilon$. Independently let
$n$ exponential gaps be conditioned to have total $n\tau$,
$\tau\in[\tau_-,\tau_+]$. For all sufficiently large $n$,
\[
 \mathbb P\left\{
 \max_{j\le n}|T_j-j\tau|\le a\sqrt n,
 \quad\max_{c,j}|N_c(j)-j\pi_c|\le a\sqrt n
 \right\}\ge c_a>0.
\]
The threshold depends on the fixed data and $a$.
\end{lemma}
\begin{proof}
First, for both exact bridges,
\begin{equation}
 \mathbb P\{\max_j|T_j-j\tau|>u\}
 +\sum_c\mathbb P\{\max_j|N_c(j)-j\pi_c|>u\}
 \le C\exp\{-c\min(u^2/n,u)\}.\label{bridge:maximal}
\end{equation}
To see this, expose at most $3n/4$ coordinates. Their likelihood
relative to independent reference coordinates is the product of
\[
 \frac{f_{n-m,\tau}(n\tau-\sum_I E_j)}{f_{n,\tau}(n\tau)},
 \qquad
 \frac{\mathbb P\{\operatorname{Mult}(n-m,\pi)=v-N_I\}}
      {\mathbb P\{\operatorname{Mult}(n,\pi)=v\}}.
\]
Gamma and multinomial maximum-density and central-density Stirling
bounds make both ratios bounded above. This is a comparison for the
whole selected set. It does not assert independence of separated
selected coordinates. Independent centered gaps and centered color
indicators have moment generating functions bounded by
$\exp(C\theta^2)$ for $|\theta|\le\theta_0$. The exponential
supermartingale maximal inequality gives the right side of
\eqref{bridge:maximal} on the first half. Reverse the second half,
using exchangeability and the zero centered complete total. This
proves \eqref{bridge:maximal}; $n=1$ is immediate.

Choose a fixed integer $J\ge2$, and split the coordinates into $J$
blocks of lengths between $n/(2J)$ and $2n/J$. In the first $J-1$
blocks, let the first $r-1$ color counts vary in integer intervals of
width $\delta\sqrt n$ about their means. Determine the remaining
color count by the block length and the last block by the full type.
All deviations are $O_r(J\delta\sqrt n)$ and all counts are
nonnegative for sufficiently large $n$. The exact conditional atom is
\[
 \frac{\prod_{b=1}^J
       \mathbb P\{\operatorname{Mult}(n_b,\pi)=v_b\}}
      {\mathbb P\{\operatorname{Mult}(n,\pi)=v\}}.
\]
Central Stirling gives a lower bound
$c_{J,\delta}n^{-(J-1)(r-1)/2}$ for each choice. There are at least
$c_{J,\delta}n^{(J-1)(r-1)/2}$ choices. Thus this event has positive
constant probability, including $r=1$.

Likewise prescribe the first $J-1$ block gap totals in real intervals
of width $\delta\sqrt n$ about $n_b\tau$, with the last determined by
the full total. Their density is
$\prod_b f_{n_b,\tau}(A_b)/f_{n,\tau}(n\tau)$, bounded below by
$c_{J,\delta}n^{-(J-1)/2}$ on a box of volume
$\asymp n^{(J-1)/2}$. Its probability is positive, and the gap and
color channels remain independent.

After $J$ is fixed, choose $\delta$ small enough that all cumulative
boundary deviations are at most $a\sqrt n/4$. Conditional on these
block totals, the blocks are independent exact bridges with compact
new parameters. By \eqref{bridge:maximal}, their combined failure
probability for tubes of radius $a\sqrt n/2$ about their own linear
interpolants is at most $CJ\exp(-ca^2J)$ for sufficiently large $n$.
Choose $J$ to make this at most $1/2$, and then choose $\delta$ and
the length threshold. Boundary and internal deviations sum to less
than $a\sqrt n$. All choices precede $n$.
\end{proof}

The gap and color conclusions can have different lengths: apply their
separate estimates at $n+1$ and $n$, respectively. In particular the
lemma applies to a bridge with $n$ marks and all $n+1$ gaps. Its large
length threshold is essential for an arbitrarily narrow tube.

\subsection{Positive drift under exact conditioning}

\begin{lemma}[Ordinary positive-drift continuation]
\label{bridge:continuation}
Fix compact positive ranges for two marks $w_0,w_1$ (which may be equal),
a speed $c$, and $\tau=P/(n+1)$, with $k/n$ in a fixed compact
subinterval of $(0,1)$. On $[0,P]$ take $n$ uniform points and,
independently, a uniform $k$-subset bearing mark $w_1$; the others bear
$w_0$. For $\varepsilon\in\{-1,1\}$ put
\[
 Z(t)=x+\varepsilon\{w_0N_0(t)+w_1N_1(t)-ct\},\qquad
 v=\varepsilon\left(\frac{(n-k)w_0+kw_1}{P}-c\right).
\]
Suppose $0<v\le v_{\max}$. For every fixed $a>0$ there are
$R(a)<\infty$ and $b(a)>0$ such that
\[
 x\ge a/v,\quad Pv^2\ge R(a)
 \quad\Longrightarrow\quad
 \mathbb P\{Z(t)\ge0\ (0\le t\le P)\}\ge b(a).
\]
This probability averages both the exact clock and exact type.
\end{lemma}
\begin{proof}
Write $n_0=n-k$, $n_1=k$. The law is equivalently that of independent
uniform positions of each color, with respective counts $n_i$. Use
independent Poisson processes of rates $\lambda_i=n_i/P$ as a reference
law $Q$, and call their law conditioned on the two totals $B$. On the
first half, the exact density is
\begin{equation}
 \frac{dB}{dQ}
 =\prod_{i=0}^1
 \frac{p_{n_i/2}(n_i-N_i(P/2))}{p_{n_i}(n_i)},
 \qquad p_\mu(r)=e^{-\mu}\mu^r/r!.
 \label{bridge:poisson-half-density}
\end{equation}
It is at most a fixed $C_{\rm den}$ everywhere. On
$\mathcal C_K=\{|N_i(P/2)-n_i/2|\le K\sqrt{n_i},\ i=0,1\}$
it is at least $c_{\rm den}(K)>0$ once $n_i\ge N_0(K)$.
The upper follows from Poisson maximum atoms, and the lower from
central Stirling, for example first taking $n_i\ge16K^2$. The same
statements apply to the reversed half.

If $w_+$ bounds the marks above and $\tau_-$ bounds the gap mean
below, put
\[
 C_{\exp}=\frac{w_+^2e^{w_+}}{2\tau_-},\qquad
 \kappa=\min\{v_{\max}^{-1},(2C_{\exp})^{-1}\}.
\]
For $0\le\theta\le1$, the Taylor remainder gives in both orientations
\[
 \varepsilon c\theta+
 \sum_i\lambda_i(e^{-\varepsilon\theta w_i}-1)
 \le-v\theta+C_{\exp}\theta^2.
\]
With $\theta=\kappa v$, the process $e^{-\theta Z(t)}$ is a
nonnegative supermartingale. At the first exit from $(0,\infty)$
its value is at least one; the possible negative overshoot is bounded
by $w_+$. Stopping at the exit or a fixed time and passing to the
limit proves that the exit probability under $Q$ is at most
$e^{-\theta x}\le e^{-\kappa a}$. Thus first-half survival has
$Q$-probability at least $p_a=1-e^{-\kappa a}>0$.
Choose $K\ge\max(1,\sqrt{2/p_a})$. Chebyshev gives
$Q(\mathcal C_K^c)\le K^{-2}\le p_a/2$, so
\[
 B(\text{first-half survival})
 \ge c_{\rm den}(K)p_a/2=:b_{\rm left}(a)>0.
\]

For $0\le u\le P/2$, let $N_i^{\rm rev}(u)$ count points in the
last interval of length $u$, with the strict-open convention at the
moving left endpoint. The exact identity is
\[
 Z(P-u)=x+v(P-u)-\varepsilon M^{\rm rev}(u),\qquad
 M^{\rm rev}(u)=\sum_iw_i(N_i^{\rm rev}(u)-\lambda_i u).
\]
A second-half failure therefore forces
$\sup_{u\le P/2}\varepsilon M^{\rm rev}(u)\ge vP/2$.
Including one-sided limits makes this the usual right-continuous
Poisson maximal event. Its reference exponential compensator is
\[
 \Lambda_\varepsilon(\theta)
 =\sum_i\lambda_i(e^{\varepsilon\theta w_i}-1
                           -\varepsilon\theta w_i)
 \le C_{\exp}\theta^2.
\]
The exponential martingale maximal inequality and the upper density
in \eqref{bridge:poisson-half-density} give
\[
 B(\text{second-half failure})
 \le C_{\rm den}e^{-\kappa Pv^2/4}.
\]
Choose $R(a)$ to make this at most $b_{\rm left}(a)/2$, and large
enough that $Pv^2\ge R(a)$ forces $n_i\ge N_0(K)$. It suffices for
the latter that
$R(a)/v_{\max}^2\ge\tau_+[N_0(K)/\epsilon+1]$.
Subtracting gives the result. All constants are fixed before the
bridge or any retained prefix is chosen.
\end{proof}

The reference rank drift is
$\varepsilon(\mathbb EW-c\tau)=(y-x+\varepsilon c\tau)/n$.
It differs from the physical drift $(y-x)/P$ by the contribution of
the terminal unmarked gap; neither gap nor contribution is omitted.
